% ============================================================
%  sections/06_discussion.tex
% ============================================================
\section{Discussion}
\label{sec:discussion}

\subsection{The FLOPs--Latency Gap for Unstructured Pruning}

Our most important empirical finding is the \emph{non-monotonic} relationship
between sparsity and wall-clock latency for unstructured pruning.
Table~\ref{tab:pruning} reveals a surprising pattern:

\begin{itemize}[leftmargin=1.5em,itemsep=2pt]
  \item \textbf{25\%--75\% sparsity}: latency \emph{increases} by
    13--41\% relative to the baseline, despite theoretical FLOPs decreasing
    proportionally. At 50\% sparsity (0.61M vs.\ 1.22M theoretical MACs),
    measured latency increases from 0.56ms to 0.80ms.
  \item \textbf{90\% sparsity}: latency \emph{decreases} to 0.34ms
    ($1.7\times$ faster than baseline), with throughput rising from
    6{,}137 to 11{,}277 samples/sec.
\end{itemize}

This non-monotonic behaviour has a hardware-level explanation rooted in
how modern CPUs execute matrix-vector products.

\paragraph{Why 25--75\% sparsity \emph{increases} latency.}
Standard CPUs execute dense matrix-vector products via BLAS Level~2 routines
(e.g., \texttt{sgemv} in MKL/OpenBLAS) that saturate SIMD units by
processing 8--16 FP32 values per instruction.
A sparse weight matrix stored in dense format performs the \emph{same}
memory loads and stores as a dense one; the zero multiplications are
dispatched to the FPU identically to non-zero ones.
Additional overhead comes from: (a)~the forward hook that re-applies
the mask before each forward pass at 25--75\% sparsity (where masks are
still active during our profiling run), and (b)~increased cache pressure
from random sparsity patterns breaking sequential memory access.

\paragraph{Why 90\% sparsity \emph{reduces} latency.}
At extreme sparsity ($>$90\%), a different regime emerges.
With only 10\% of weights non-zero in the dominant FC1 layer
($200{,}832$ parameters $\to$ $\sim$20{,}000 non-zero),
the CPU's branch predictor and arithmetic unit pipeline begins to
exhibit \emph{zero-skipping} behaviour in the accumulation loop.
Furthermore, at this sparsity level, all non-zero weights fit within
L2 cache, eliminating DRAM latency for the hot path.
This is a \textbf{hardware-specific} effect: the same model may exhibit
different latency on different CPU microarchitectures.

\paragraph{Implication for practitioners.}
The standard claim ``unstructured pruning does not reduce inference
latency'' is an \emph{oversimplification}.
The correct statement is:
\begin{enumerate}[leftmargin=2em,topsep=2pt]
  \item At moderate sparsity (25--75\%), unstructured pruning
    \emph{increases} measured latency due to masking overhead and
    cache-access pattern degradation.
  \item At extreme sparsity ($\geq$90\%), hardware effects can produce
    latency \emph{reduction}, but this is not predictable across
    architectures and should not be relied upon for deployment.
  \item Real, reliable latency benefits require structured pruning,
    sparse-aware kernels, or specialized hardware.
\end{enumerate}


\subsection{Quantization as the Practical Default}

Dynamic INT8 PTQ achieves the best trade-off on all Pareto dimensions
simultaneously, and requires no additional training.
The accuracy cost ($<$0.1\% degradation) is negligible for most
applications.

Two caveats deserve attention:
\begin{enumerate}[leftmargin=2em]
  \item \textbf{Conv2d is not quantized.}
    PyTorch's \texttt{quantize\_dynamic} targets Linear layers only.
    For models where Conv2d operations dominate (e.g., ResNet),
    static PTQ or quantization-aware training (QAT) is needed.
  \item \textbf{Deprecation.}
    The \texttt{torch.quantization} API is deprecated in PyTorch 2.x
    and will be removed in 2.10.
    The migration path is \texttt{torchao} (PyTorch's quantization library).
\end{enumerate}

\subsection{Low-Rank Factorization: When Does it Help?}

Our results show that SVD reconstruction preserves accuracy well at
rank ratios $\geq$50\%, which aligns with the spectral theory of
over-parameterized networks: the weight matrices of trained networks
tend to have low effective rank~\citep{li2018constrained,hu2021lora}.

However, without architectural factorization, neither memory nor
latency improves. The practical value of low-rank approximation is
therefore as:
(1) a weight initialization for factorized architectures,
(2) a compression analysis tool (how compressible is a layer?), or
(3) a component in combined pipelines (prune $\to$ SVD $\to$ quantize).

\subsection{Accuracy Robustness at Extreme Sparsity}

A notable finding is that pruning to 90\% sparsity retains $>$91\%
test accuracy, despite removing 9 out of every 10 weights.
This is consistent with the lottery ticket hypothesis~\citep{frankle2019lottery}:
a small sparse subnetwork exists within the original model that
captures most of its predictive capacity.
Fine-tuning for only 5 epochs after pruning is sufficient to recover
most of the accuracy loss.

\subsection{Limitations}

\begin{enumerate}[leftmargin=2em,topsep=2pt]
  \item \textbf{Single architecture and dataset.} Conclusions are drawn
    from SmallCNN on Fashion-MNIST. Generalization to larger models
    (ResNet, ViT) and datasets (ImageNet, CIFAR-100) requires further study.
  \item \textbf{CPU-only.} GPU latency profiles differ significantly;
    INT8 benefits on GPU depend on Tensor Core availability.
  \item \textbf{No combined methods.} Pruning + quantization combined
    (a common production approach) is not evaluated.
  \item \textbf{Memory underestimation.} \texttt{tracemalloc}
    underestimates actual memory; process RSS is a better but
    coarser metric.
\end{enumerate}
